\documentclass[aspectratio=169,11pt]{beamer}

\usetheme{Rochester}
\usecolortheme{beaver}
\usefonttheme{serif}
\usepackage[utf8]{inputenc}
\usepackage[T1]{fontenc}
\usepackage[spanish,es-nodecimaldot]{babel}
\usepackage{amsmath,amssymb,amsthm}
\usepackage{booktabs}
\usepackage{array}

\setbeamertemplate{navigation symbols}{}
\setbeamertemplate{theorems}[numbered]
\newtheorem{proposicion}{Proposición}
\newtheorem{teorema}{Teorema}
\newtheorem{definicion}{Definición}
\newtheorem{observacion}{Observación}
\newcommand{\Z}{\mathbb Z}
\newcommand{\Zn}{\mathbb Z_n}
\newcommand{\U}{U(n)}
\newcommand{\modn}[1]{\pmod{#1}}

\title[Aplicaciones de los productos directos]{Aplicaciones de los productos directos}
\subtitle{Seguridad de datos, criptografía de clave pública y firmas digitales}
\author[L. Hernández]{Lester Iván Hernández Hernández\\\small 20242000637\\[2mm]
Jason Dionicio Ventura\\\small 20172031480}
\institute[UNAH--Campus Cortés]{Estructuras Algebraicas I (MM-421)\\
Universidad Nacional Autónoma de Honduras (UNAH), Campus Cortés\\
Docente: Lic. Edwin Hernández}
\date{Martes 17 de noviembre de 2026}

\begin{document}

\begin{frame}
  \titlepage
\end{frame}

\begin{frame}{Objetivo y mapa de ideas}
  \begin{itemize}
    \item Modelar cadenas de datos mediante productos internos directos de grupos cíclicos.
    \item Entender por qué la aritmética modular permite cifrar y descifrar.
    \item Ver cómo la inversión de algoritmos produce autenticidad, no solo secreto.
  \end{itemize}
  \vfill
  \begin{center}
    $(\mathbb Z_2)^n \longrightarrow$ OTP \qquad
    $U(n)\cong U(p)\oplus U(q) \longrightarrow$ RSA \qquad
    $D_A/E_A \longrightarrow$ firma digital
  \end{center}
\end{frame}

\begin{frame}{La conexión algebraica común}
  \begin{itemize}
    \item Un producto interno directo $G_1\oplus\cdots\oplus G_k$ opera componente a componente:
    \[
      (a_1,\ldots,a_k)(b_1,\ldots,b_k)=(a_1b_1,\ldots,a_kb_k).
    \]
    \item En $(\mathbb Z_2)^n$ la operación es suma módulo $2$; cada bit es una coordenada.
    \item En $U(n)$ la operación es multiplicación módulo $n$. Si $n=pq$ con primos distintos,
    \[
      U(n)\cong U(p)\oplus U(q)\cong \mathbb Z_{p-1}\oplus\mathbb Z_{q-1}
    \]
    cuando se identifican los grupos cíclicos apropiados.
  \end{itemize}
\end{frame}

\section{Data Security}

\begin{frame}{Data Security: el problema motivador}
  \begin{itemize}
    \item Una transacción necesita viajar por un canal que puede ser observado.
    \item Si enviamos el mensaje directamente, un interceptor obtiene la información.
    \item Un cifrado simple también puede revelar patrones si la clave se reutiliza.
    \item Solución algebraica: sumar una cadena aleatoria del mismo tamaño y usarla una sola vez.
  \end{itemize}
\end{frame}

\begin{frame}{Definición formal: libreta de un solo uso}
  Sea $M\in(\mathbb Z_2)^n$ el mensaje y sea $K\in(\mathbb Z_2)^n$ una clave aleatoria.
  \[
    C=M+K,\qquad \text{Descifrado}(C)=C+K.
  \]
  \begin{definicion}[Cifrado OTP]
    El cifrado de libreta de un solo uso es la transformación
    $E_K:(\mathbb Z_2)^n\to(\mathbb Z_2)^n$, $E_K(M)=M+K$,
    donde $K$ se genera uniformemente y no se vuelve a utilizar. La abreviatura
    \textbf{OTP} viene del inglés \textit{One-Time Pad}.
  \end{definicion}
  La suma es componente a componente; por tanto, es exactamente la operación del producto directo $(\mathbb Z_2)^n$.
\end{frame}

\begin{frame}{Ejemplo resuelto: cifrar}
  Mensaje de 10 bits y clave aleatoria:
  \[
    M=1101010110,\qquad K=0110011001.
  \]
  Sumamos cada coordenada módulo $2$:
  \[
  \begin{array}{lrrrrrrrrrr}
  M:&1&1&0&1&0&1&0&1&1&0\\
  K:&0&1&1&0&0&1&1&0&0&1\\ \midrule
  C:&1&0&1&1&0&0&1&1&1&1
  \end{array}
  \]
  Por tanto, se transmite $C=1011001111$, no $M$.
\end{frame}

\begin{frame}{Ejemplo resuelto: descifrar y justificar}
  El receptor conoce $K$ y calcula:
  \[
  C+K=1011001111+0110011001=1101010110=M.
  \]
  \begin{proposicion}
    Para todo $M,K\in(\mathbb Z_2)^n$, $E_K(E_K(M))=M$.
  \end{proposicion}
  \begin{proof}
    $E_K(E_K(M))=(M+K)+K=M+(K+K)=M+0=M$, porque cada coordenada satisface $1+1=0$ en $\mathbb Z_2$.
  \end{proof}
\end{frame}

\begin{frame}{¿Por qué la clave debe usarse una sola vez?}
  Si se reutilizan $K$ para dos mensajes,
  \[
    C_1=M_1+K,\qquad C_2=M_2+K,
  \]
  entonces un interceptor puede calcular
  \[
    C_1+C_2=M_1+M_2.
  \]
  La clave desaparece y aparecen relaciones entre los mensajes. La seguridad perfecta del OTP requiere:
  \begin{enumerate}
    \item clave uniformemente aleatoria;
    \item longitud de clave igual a la del mensaje;
    \item distribución y reutilización imposibles de predecir.
  \end{enumerate}
\end{frame}

\section{Public Key Cryptography}

\begin{frame}{Criptografía de clave pública: el problema}
  En un sistema simétrico, emisor y receptor deben compartir una clave secreta antes de comunicarse.
  \begin{block}{Pregunta central}
    ¿Cómo puede un receptor publicar la manera de cifrar sin publicar la manera de descifrar?
  \end{block}
  RSA (Rivest, Shamir y Adleman) aprovecha la asimetría computacional:
  multiplicar primos grandes es fácil, factorizar su producto es difícil.
\end{frame}

\begin{frame}{Construcción RSA}
  El receptor elige primos distintos $p,q$ y define $n=pq$ y
  $m=\operatorname{mcm}(p-1,q-1)$, donde \textbf{mcm} significa mínimo común múltiplo.
  \begin{enumerate}
    \item Escoge $e$ con $\gcd(e,m)=1$.
    \item Calcula $d$ tal que $ed\equiv1\modn{m}$.
    \item Publica $(n,e)$ y conserva $p,q,d$.
  \end{enumerate}
  Para un bloque $M\in U(n)$:
  \[
    R=M^e\modn{n},\qquad M=R^d\modn{n}.
  \]
  El grupo $U(n)$ explica por qué la exponenciación es invertible.
\end{frame}

\begin{frame}{Ejemplo RSA: generación de claves}
  Tomemos un ejemplo pequeño, solo para poder hacer las cuentas:
  \[
    p=11,\quad q=17,\quad n=187,\quad m=\operatorname{mcm}(10,16)=80.
  \]
  Elegimos $e=3$, pues $\gcd(3,80)=1$. El inverso de $3$ en $U(80)$ es
  \[
    d=27,\qquad ed=3\cdot27=81\equiv1\modn{80}.
  \]
  \begin{block}{Claves}
    Pública: $(187,3)$. \qquad Privada: $27$ (y la factorización $187=11\cdot17$).
  \end{block}
\end{frame}

\begin{frame}{Ejemplo RSA: cifrado paso a paso}
  Sea $M=42$. Como $\gcd(42,187)=1$, $M\in U(187)$.
  \[
    R=42^3\modn{187}.
  \]
  Calculamos:
  \[
    42^2=1764\equiv81\modn{187},\qquad
    42^3\equiv81\cdot42=3402\equiv36\modn{187}.
  \]
  El mensaje transmitido es, entonces, $R=36$.
\end{frame}

\begin{frame}{Ejemplo RSA: descifrado por exponenciación rápida}
  Calculamos $36^{27}\modn{187}$ usando $27=16+8+2+1$:
  \[
  \begin{array}{c|ccccc}
  \text{potencia}&36^1&36^2&36^4&36^8&36^{16}\\ \midrule
  \text{módulo }187&36&174&169&137&69
  \end{array}
  \]
  Entonces
  \[
    36^{27}\equiv36^{16}36^8 36^2 36
    \equiv69\cdot137\cdot174\cdot36\equiv42\modn{187}.
  \]
  Recuperamos exactamente $M=42$.
\end{frame}

\begin{frame}{Demostración: por qué RSA funciona}
  Como $n=pq$ y $M\in U(n)$, el teorema chino del resto da
  \[
    U(n)\cong U(p)\oplus U(q)\cong\mathbb Z_{p-1}\oplus\mathbb Z_{q-1}.
  \]
  Si $m=\operatorname{mcm}(p-1,q-1)$, entonces para $x\in U(n)$:
  \[
    x^m\leftrightarrow\bigl(x_1^m,x_2^m\bigr)=(1,1),
  \]
  porque $p-1\mid m$ y $q-1\mid m$.
  Por tanto, $x^m=1$ en $U(n)$.
\end{frame}

\begin{frame}{Demostración RSA: inversión de exponentes}
  Puesto que $ed\equiv1\modn{m}$, existe $k\in\Z$ tal que $ed=1+km$.
  Para $M\in U(n)$:
  \[
    R^d=(M^e)^d=M^{ed}=M^{1+km}=M(M^m)^k=M\cdot1^k=M.
  \]
  La demostración usa exactamente la estructura de grupo y el producto directo.
  \begin{alertblock}{Alcance}
    Esta exposición cita como hecho criptográfico que factorizar productos de primos grandes es difícil; no es una consecuencia de la teoría de grupos.
  \end{alertblock}
\end{frame}

\section{Digital Signatures}

\begin{frame}{Firmas digitales: autenticidad}
  Cifrar garantiza secreto, pero no responde por sí solo: «¿quién envió este mensaje?»
  \begin{itemize}
    \item Un banco necesita distinguir una orden auténtica de una orden falsificada.
    \item La firma debe ser verificable públicamente.
    \item Solo la persona firmante debe poder producir una firma válida.
  \end{itemize}
\end{frame}

\begin{frame}{Definición y protocolo}
  Para A, $E_A$ es pública y $D_A$ es privada; para B, $E_B$ es pública y $D_B$ es privada.
  A envía el mensaje firmado y confidencial:
  \[
    T=E_B(D_A(M)).
  \]
  B aplica $D_B$ y luego $E_A$:
  \[
    E_A(D_B(T))=E_A(D_B(E_B(D_A(M))))=E_A(D_A(M))=M.
  \]
  El orden importa: A prueba autoría con $D_A$ y B obtiene secreto con $E_B$.
\end{frame}

\begin{frame}{Ejemplo: firma RSA de un resumen}
  A usa $n_A=187$, $e_A=3$, $d_A=27$; B usa $n_B=221$, $e_B=5$, $d_B=29$.
  Para un mensaje, ambos calculan el resumen $h=50$.
  \[
    s=h^{d_A}=50^{27}\modn{187}=118.
  \]
  A protege la firma para B:
  \[
    T=s^{e_B}=118^5\modn{221}=118.
  \]
  B descifra con $d_B$: $T^{d_B}\modn{221}=118$.
\end{frame}

\begin{frame}{Verificación y demostración}
  B verifica la firma aplicando la clave pública de A:
  \[
    s^{e_A}=118^3\modn{187}=50=h.
  \]
  \begin{proposicion}
    Si $e_A d_A\equiv1\modn{m_A}$, entonces una firma $s=h^{d_A}$ satisface $s^{e_A}=h$ para $h\in U(n_A)$.
  \end{proposicion}
  \begin{proof}
    $s^{e_A}=(h^{d_A})^{e_A}=h^{d_Ae_A}=h^{1+km_A}=h(h^{m_A})^k=h$,
    por el mismo argumento de RSA.
  \end{proof}
\end{frame}

\begin{frame}{Qué garantiza y qué no garantiza una firma}
  \begin{columns}[T]
    \column{.48\textwidth}
    \textbf{Sí garantiza, bajo las hipótesis del esquema:}
    \begin{itemize}
      \item integridad del mensaje resumido;
      \item evidencia de posesión de la clave privada;
      \item verificación pública.
    \end{itemize}
    \column{.48\textwidth}
    \textbf{No resuelve por sí sola:}
    \begin{itemize}
      \item robo de la clave privada;
      \item colisiones o debilidades del resumen;
      \item distribución inicial de claves.
    \end{itemize}
  \end{columns}
\end{frame}

\section{Problema de aplicación computacional}

\begin{frame}{Problema de aplicación computacional}
  Una institución recibe $10\,000$ mensajes. Cada mensaje contiene un bloque $M_i<n$, una firma $s_i$ y la clave pública RSA $(n,e)$.
  \begin{block}{Entrada y salida}
    Entrada: un entero $N=10\,000$, los enteros $n=3233$, $e=17$ y $N$ parejas $(M_i,s_i)$, con $0<M_i<n$.\\
    Salida: los índices $i$ para los que $s_i^e\not\equiv M_i\modn{n}$ y el total de firmas inválidas.
  \end{block}
  En este ejercicio cada bloque $M_i$ se firma directamente, sin una función resumen externa, para que todo pueda resolverse en un único programa.
\end{frame}

\begin{frame}{Parámetros del ejercicio}
  \[
    n=3233=61\cdot53,\qquad e=17.
  \]
  \begin{itemize}
    \item Estos valores corresponden a una clave pública RSA de demostración.
    \item El programa solo necesita $n$ y $e$; no necesita conocer los factores $61$ y $53$ ni la clave privada.
    \item Se procesan $10\,000$ parejas, suficientes para que la automatización sea necesaria sin requerir bibliotecas criptográficas externas.
  \end{itemize}
\end{frame}

\begin{frame}{Pseudocódigo del verificador}
  \begin{enumerate}
    \item Leer $N$, $n$, $e$.
    \item Repetir $N$ veces: leer $(M,s)$ y calcular $v=\operatorname{potencia\_modular}(s,e,n)$.
    \item Si $v\ne M$, guardar el índice de la pareja.
    \item Imprimir los índices guardados y su cantidad.
  \end{enumerate}
  \vfill
  La función \texttt{potencia\_modular} puede implementarse con el método de cuadrados sucesivos; en Python o C++ cabe en un mismo programa pequeño.
\end{frame}

\begin{frame}{Tema aplicado y algoritmo}
  \textbf{Tema:} firmas digitales; la igualdad que se comprueba es la verificación RSA del bloque firmado.
  \begin{enumerate}
    \item Leer cada pareja $(M_i,s_i)$.
    \item Calcular $v_i=s_i^{17}\modn{3233}$ mediante exponenciación binaria.
    \item Marcar $i$ como inválido si $v_i\ne M_i$.
  \end{enumerate}
  No se incluye código: basta implementar multiplicación modular segura y exponenciación por cuadrados.
\end{frame}

\begin{frame}{Corrección y complejidad aproximada}
  \begin{itemize}
    \item Si una firma fue producida por la clave privada correspondiente, la demostración RSA da $s_i^{17}\equiv M_i\modn{3233}$; el algoritmo acepta.
    \item Si la igualdad falla, el par no satisface la condición de verificación; el algoritmo lo rechaza.
    \item Para $N$ mensajes y exponentes de $b$ bits, la exponenciación binaria usa $O(b)$ multiplicaciones modulares por mensaje: $O(Nb)$ multiplicaciones, ignorando el costo de multiplicar enteros.
    \item Aquí $N=10\,000$ y $17$ tiene solo cinco bits: el programa es corto, pero revisar todas las firmas a mano no es razonable.
  \end{itemize}
\end{frame}

\begin{frame}{Bibliografía}
  \begin{thebibliography}{9}
    \bibitem{gallian}
    Joseph A. Gallian, \textit{Contemporary Abstract Algebra}, 9.ª ed., Cengage Learning, 2015, capítulo 8, sección ``Applications'', pp. 162--167.
    \bibitem{rsa}
    R. L. Rivest, A. Shamir y L. Adleman, ``A Method for Obtaining Digital Signatures and Public-Key Cryptosystems'', \textit{Communications of the ACM}, 21(2), 1978, pp. 120--126.
  \end{thebibliography}
  \vfill
  \centering\small
  Los ejemplos numéricos de esta presentación fueron construidos para la exposición y no reproducen los ejemplos del libro.
\end{frame}

\end{document}
